% Generated draft: requires mathematical and source review.
\documentclass[11pt]{article}
\usepackage[margin=25mm]{geometry}
\usepackage{fontspec}
\setmainfont{Latin Modern Roman}
\usepackage{amsmath,amssymb,mathtools,unicode-math}
\setmathfont{Latin Modern Math}
\usepackage{enumitem,needspace,xurl,hyperref}
\setlength{\parindent}{0pt}
\setlength{\parskip}{5pt}
\setlength{\emergencystretch}{3em}
\newcommand{\N}{\mathbb{N}}
\newcommand{\Z}{\mathbb{Z}}
\newcommand{\Q}{\mathbb{Q}}
\newcommand{\R}{\mathbb{R}}
\newcommand{\C}{\mathbb{C}}
\newcommand{\F}{\mathbb{F}}
\newcommand{\A}{\mathbb{A}}
\newcommand{\PP}{\mathbb{P}}
\newcommand{\E}{\mathbb{E}}
\newcommand{\eps}{\varepsilon}
\newcommand{\dd}{\,\mathrm{d}}
\begin{document}
\section*{UnsolvedMath ID: 30001759}
\textbf{Source problem code:} catalogue entry
\subsection{Definitions and mathematical statement}
Let \(N\geq 0\) and let \(Y\subset\mathbb{CP}^{N}\) be a complex projective variety. Write
\[
Y_{\mathrm{reg}}=Y\setminus\operatorname{Sing}Y,
\qquad g=g_{\mathrm{FS}}|_{Y_{\mathrm{reg}}},
\]
where \(g_{\mathrm{FS}}\) is the Fubini--Study metric. For a coefficient choice \(\mathbb K\in\{\R,\C\}\), let \(L^{2}\Omega^{k}(Y_{\mathrm{reg}};\mathbb K)\) denote the measurable differential \(k\)-forms with finite squared norm integrated against the volume measure of \(g\). The source does not specify the coefficient choice.

The exterior derivative initially acts on compactly supported smooth forms. On this possibly incomplete manifold, two relevant closed realizations are its minimal extension \(d_{\min}\), obtained by graph closure of that operator, and its maximal extension \(d_{\max}\), whose domain consists of \(L^{2}\)-forms whose distributional exterior derivative is also \(L^{2}\). For a specified closed realization \(d_D\) forming a differential complex, its unreduced cohomology is
\[
H^{k}_{(2),D}(Y_{\mathrm{reg}};\mathbb K)
=
\frac{\ker(d_D^{k})}{\operatorname{im}(d_D^{k-1})},
\]
whereas its reduced cohomology is
\[
\overline H^{k}_{(2),D}(Y_{\mathrm{reg}};\mathbb K)
=
\frac{\ker(d_D^{k})}{\overline{\operatorname{im}(d_D^{k-1})}}.
\]
The closure is in the \(L^{2}\) norm. In degree zero the denominator is zero. The symbol \(D\) records missing analytic data; the question below does not quantify over all realizations.

Write \(IH^{k}_{\bar m}(Y;\mathbb K)\) for intersection cohomology with middle perversity. On a pure-dimensional stratified pseudomanifold, the lower middle perversity has value \(\bar m(c)=\lfloor(c-2)/2\rfloor\) at real codimension \(c\geq2\); lower and upper middle perversities agree on even codimensions. The source does not specify conventions for varieties that are not pure-dimensional.

\textbf{Source question, with unresolved conventions displayed.} For every projective variety \(Y\subset\mathbb{CP}^{N}\) and every degree \(k\geq0\), is there a natural isomorphism
\[
H^{k}_{(2),D}(Y_{\mathrm{reg}};\mathbb K)
\stackrel{?}{\cong}
IH^{k}_{\bar m}(Y;\mathbb K)
\]
when \(Y_{\mathrm{reg}}\) carries the restricted Fubini--Study metric? If reduced cohomology was intended, the left-hand side must instead be \(\overline H^{k}_{(2),D}\). The supplied statement does not determine \(D\), reduced versus unreduced cohomology, coefficients, the intended dimension conventions, or the meaning of naturality. Consequently this displays the intended comparison but does not resolve it into a unique mathematical assertion. No additional analytic hypothesis is imposed here.
\subsection{Short English statement}
The question asks whether square-integrable de Rham cohomology on the smooth part of every complex projective variety, equipped with the induced Fubini--Study metric, naturally recovers its middle-perversity intersection cohomology in every degree. The differential domain, reduction convention, coefficients, and some geometric conventions remain unspecified. The supplied partially solved label is dated curation, not an independently verified conclusion.
\subsection{Sources}
\begin{itemize}
\item Local supplied UnsolvedMath registry record: permanent ID \(30001759\), source problem code \texttt{catalogue entry}, in the collection \emph{Oberwolfach Reports Open Problems}. The supplied record gives proposed year 2011 and names no individual proposer.
\item \emph{Geometric Methods of Complex Analysis} (2011), Oberwolfach Reports, pp.\ 1047--1116, DOI 10.4171/owr/2011/19. Supplied workshop-report provenance: \url{https://doi.org/10.4171/owr/2011/19}. No precise problem page or passage is supplied.
\item Supplied literature reference, attributed to Cheeger--Goresky--MacPherson: \emph{L2 cohomology and intersection homology}, \url{https://www.math.ias.edu/~goresky/pdf/CGM.pdf}.
\item Supplied literature reference labelled ``recent analysis of the CGM conjecture'': \url{https://arxiv.org/abs/2508.01679}. The record supplies no authors or full bibliographic title for this reference.
\item Literature assessment dated \texttt{2026-08-21}: the supplied curation labels the problem partially solved and states that the Cheeger--Goresky--MacPherson \(L^{2}\)/intersection-cohomology identification is proved in important projective and isolated-singularity cases, while the unrestricted extracted formulation remains subject to analytic hypotheses. This is supplied dated curation, not independently verified; its time-sensitive status and cited claims require independent checking.
\item The registry labels the statement \texttt{corrected\_verified} and reports a manual semantic, typographic, and LaTeX review of the cleanup. These are supplied editorial assertions, not mathematical verification of this draft or of the conjecture.
\end{itemize}
\end{document}
